% Pursuing Stacks (cotes 134-2 to 134-8): archive cote 134-4, site page 20 (PDF page 21), typescript page 108.
% typescript. Draft compared: ps3.tex:978-1039.
% AI-transcribed by Claude Opus 5.5 (claude-opus-5-5, adaptive thinking, effort high), Message Batches API, batch ps1 (msgbatch_01Rv5hCaNTtUt4JHx1yqoHPL), 2026-11-12; not yet checked by eye.
\scanpage{134-4/20}
\mspage{108}
tible,  contrarily to what one would expect from the intuitive meaning of
an ``interval''.  In those two examples,  the condition of getting two
``disjoint'' sections of the aspheric I over e  looks kind of trivial - \xout{kix}
kind of unrelated to the question of whether or not I is actually contrac-%
tible, \inserted{and} not only aspheric.  What comes to mind here is to look for \deleted{strict} \inserted{contractible}
\uline{homotopy intervals} ,  namely homotopy intervals \inserted{$(I,\delta_0,\delta_1)$} (including condition b)
that Ker($\hand{$\delta$}_0,\hand{$\delta$}_1$) = $\emptyset_M$ )  \deleted{which are} \inserted{such that $I$ is} not only aspheric \inserted[typed]{over e}, but  \deleted{which are} \inserted{even} con-%
tractible.  Existence of \deleted{strict} \inserted{contractible} homotopy intervals seems a priori a lot
stronger than just existence of homotopy intervals,  but in the cases I'
ve met so far,  the two apparently coincide.  As a matter of fact,  the
\xout{most} first choice of homotopy interval that comes to mind,  in all those
examples,  \uline{is} indeed a \deleted{strict} \inserted{contractible} one  -  so much so that till the notes today
I was somewhat confused on this matter,  and was under the tacit impressi-%
on that homotopy intervals are all contractible,  and trivially so ! But
probably,  as usual \deleted{under this} \inserted{with most} evidently false impression\unsure{\hand{s,}|s,} \inserted{underneath,} something cor-%
rect \xout{and} \inserted[typed]{should exist,} worth \xout{whix} being explicitly stated.

First thing that comes to mind is here to look at the simplest case
of an aspheric interval \xout{which} \uline{I} which is contractible  -  namely when
\uline{I} is even \uline{I}-contractible,  and more specifically still, \xout{xx} \inserted[typed]{when} there exists
a\xout{nhomotopy} \uline{I}-homotopy (an ``elementary'',  not a ``composed'' one !) \xout{brxn}
from id$_I$ to one of the two constant maps of I into itself,  defined by the
two sections $\hand{$\delta$}_0$,$\hand{$\delta$}_1$  - say by $\hand{$\delta$}_1$ , to be specific.  Such a homotopy is a
map
\[ h\colon IxI \to I \]
in M,  having the two properties expressed symbolically by
\[ h(e_0,x) = x \quad , \quad h(e_1,x) = e_1 \quad , \]
where x may be viewed as any ``point'' of I with ``values'' in an arbitrary
parameter object T of M,  i.e.  $x : T \to I$ ,  and $e_0$ , $e_1$ are the constant
maps $T \to I$ defined by $\hand{$\delta$}_0,\hand{$\delta$}_1$.  If we view h as a composition law on I,
these relations just mean that $e_0$ acts as a left unit,  and $e_1$ as a left
zero element.  Such composition law in an interval I,  a would-be homotopy
interval as a matter of fact,  has been repetitively used before,  and
this use systematized in the ``comparison lemma for homotopy intervals''
(page 60);  \over{T}{t}here we found that if I is an object in M with such a composi-%
tion law,  and if \xout{Jxix} \inserted[typed]{there exists} an \uline{aspheric} interval \uline{J} and a morphism of intervals
from \uline{J} into \uline{I} ,  then I is equally aspheric. \unsure{\over{t}{T}|\hand{T}}his we would now see as
coming from the fact that I being \uline{I}-contractible is  \uline{J}-contractible, and
a fortiori (as \uline{J} is aspheric over e) aspheric over e. \hand{In any case,}
\hand{we see that} when I is endowed with a composition law as above, then
I is aspheric over e iff it is contractible,  and equivalently iff \hand{$\delta_0$} and
\hand{$\delta_1$} are \inserted{W-}homotopic. \margin{20}
